The rapid adoption of Large Language Models (LLMs) in service operations—from customer support chatbots to intelligent search systems—has created unprecedented computational bottlenecks that directly impact service quality and operational costs. We address a critical resource allocation challenge: optimizing LLM inference under memory constraints, where the inference process generates responses by incrementally expanding a memory-intensive Key-Value (KV) cache. Daily inference costs can exceed \$700,000 for systems like ChatGPT, making efficient scheduling essential for sustainable service delivery. We formulate LLM inference optimization as a multi-stage online scheduling problem with stochastic arrivals and dynamic resource consumption. Traditional scheduling algorithms fail because the KV cache grows unpredictably during inference, rendering conventional approaches like Shortest Job First ineffective. To overcome this challenge, we develop a fluid dynamics approximation that provides a tractable performance benchmark and guides algorithm design. Building on this theoretical foundation, we introduce two novel scheduling algorithms. First, the Waiting for Accumulated Inference Threshold (WAIT) algorithm optimizes scheduling when output lengths are known, using dynamically maintained thresholds to determine processing order. Second, for the more realistic scenario where output lengths are unknown at arrival, we extend our approach through the Nested WAIT algorithm, which adaptively learns prompt characteristics during processing through a hierarchical multi-segment framework. We establish theoretical near-optimality guarantees for both algorithms under heavy traffic conditions, balancing throughput, latency, and Time to First Token (TTFT). Extensive numerical experiments using the Llama-7B model on A100 GPUs, with both synthetic and real-world datasets, demonstrate superior throughput and reduced latency compared to widely adopted baselines (vLLM and Sarathi). This research bridges operations research and machine learning, providing service system managers with a theoretically grounded, practically implementable framework for efficient large-scale AI deployment.